\documentclass[11pt]{article}
\usepackage[T1]{fontenc}
\usepackage{lmodern,amsmath,amssymb,amsthm}
\usepackage[margin=0.95in]{geometry}
\usepackage[colorlinks=true,allcolors=blue]{hyperref}
\usepackage{microtype}
\newtheorem{theorem}{Theorem}
\newtheorem{lemma}{Lemma}
\newcommand{\area}{\mathcal A}
\newcommand{\CC}{\mathbb C}
\newcommand{\RR}{\mathbb R}
\newcommand{\pp}{\pi}
\DeclareMathOperator{\diag}{diag}
\title{A flag-curve proof of an outer-pedal area invariant}
\author{Alec Kriebel\\
\small Independent researcher\\
\small \href{https://orcid.org/0009-0001-9320-500X}{ORCID: 0009-0001-9320-500X}}
\date{16 October 2026}
\begin{document}
\maketitle
\begin{abstract}
For a noncircular ellipse and a proper nested confocal elliptical caustic with
least billiard period \(N\equiv2\pmod4\), the signed area of the outer tangent
polygon times that of its pedal polygon is independent of phase, for every
fixed pedal point. We give a flag-curve proof: on the quotient by the
Poncelet translation, the original area and the pedal area have at most the
same two simple poles. Riemann--Roch and reversal make them proportional;
the credited original/outer area-product law finishes the argument.
A short verification of that law in the required parity is included.
The theorem, the six-period formula, and the repeated-triangle qualification
also occur in Ferudun's contemporary preprint, which is explicitly credited.
This note offers an alternative presentation, without claiming theorem
novelty, first disclosure, or verified independence. Exact and numerical
controls accompany the analytic proof.
\end{abstract}

\section{Statement and relation to prior work}
Fix
\[
 E:\frac{x^2}{a^2}+\frac{y^2}{b^2}=1,\qquad
 C:\frac{x^2}{\alpha^2}+\frac{y^2}{\beta^2}=1,\qquad
 a>b>0,\quad \alpha^2=a^2-\lambda,\quad\beta^2=b^2-\lambda,\quad
 0<\lambda<b^2 .
\]
Suppose the oriented Poncelet billiard map has \emph{least} period
\(N=2m\), where \(m\geq3\) is odd. Write \(P_i\) for its boundary vertices
in traversal order, \(\ell_i\) for the tangent to \(E\) at \(P_i\), and
\(V_i=\ell_i\cap\ell_{i+1}\), with indices modulo \(N\).
For a fixed finite \(M\in\RR^2\), let \(Q_i(M)\) be the perpendicular foot
from \(M\) to the entire line \(\ell_i\). Define
\[
 \area(Z)=\frac12\sum_{i=0}^{N-1}\det(Z_i,Z_{i+1}),\qquad
 A=\area(P),\quad A'=\area(V),\quad B_M=\area(Q(M)).
\]
Thus the feet are on the sides of the \emph{outer tangent polygon},
not on the original billiard chords. These are signed areas; for a star
polygon they are not the areas of its filled faces.

\begin{theorem}\label{thm}
For the fixed conic pair and every fixed \(M\), \(A'B_M\) is independent
of the initial billiard phase. All admitted coprime star windings and
their reversals are included, and zero or negative \(B_M\) is allowed.
\end{theorem}

This is the least-period interpretation of the invariant \(k_{303,a}\)
tabulated by Reznik--Garcia--Koiller \cite[Section 3.4, Table 4]{rgk}
and in their arXiv edition \cite[Table 4, p.~6]{rgkarxiv}.
The rows specify \(N\equiv2\pmod4\) and every pedal point, but do not
literally print \emph{least period}. Their historical proof question marks
do not establish current openness. The repetition caveat below explains
why that convention matters.

Ferudun \cite[Theorem 1]{ferudun10}, retained in \cite{ferudun11}, states
the same theorem with the same signed, least-period and arbitrary-point
scope. That work also gives radial dependence on \(M\), the general
six-period constant, and a repeated-triangle countercontrol. Its Jacobi
coordinate proof uses complementary reversal zeros and pole cancellation.
The present quotient-divisor formulation differs in presentation; both
use central symmetry, isotropic projection poles, parallel residues, and
reversal. Neither the theorem nor those mechanisms are claimed new here.
Even the proportionality \(B_M=c_M A\) below follows from the product
theorem together with the nonzero real original/outer area product, so
rewriting it as a proportionality does not establish theorem novelty.

The flag-curve framework and the original/outer area-product law are
credited to the established Poncelet theory used by Ch\'avez-Caliz
\cite[Section 2 and Theorems 3, 6]{chavez}. We supply the required local
pole checks explicitly, and verify the latter law in the present parity
in Section~\ref{aa}. The proof is analytic, not an inference from sampled
periods.

\section{The flag curve and the two original-area poles}
Let \(X\) be the smooth complex projective flag curve of pairs \((P,L)\)
with \(P\in E\) and \(L\) a tangent to \(C\) through \(P\).
Its projection \(\pp:X\to E\) has degree two and branches at \(E\cap C\).
Putting \(c^2=a^2-b^2\), those intersections satisfy
\[
 x^2=\frac{a^2\alpha^2}{c^2},\qquad
 y^2=-\frac{b^2\beta^2}{c^2}.
\]
They are four distinct finite transverse intersections: all coordinates
are nonzero, and
\(1/(a^2\beta^2)-1/(b^2\alpha^2)
=\lambda c^2/(a^2b^2\alpha^2\beta^2)\ne0\).
The double cover is smooth and connected, so Riemann--Hurwitz gives
genus one.

Let \(\sigma\) exchange the two \(E\)-intersections of \(L\), and let
\(\tau\) exchange the two \(C\)-tangents through \(P\).
Use right-to-left composition and put \(T=\tau\sigma\). These are the
standard Poncelet involutions \cite[Section 2]{chavez}; their composition
is a translation of exact order \(N\). Equivalently, on a genus-one curve
each branched deck involution acts as minus the identity on holomorphic
differentials, so their composition is a translation; closure and least
period make its order exactly \(N\). We have
\[
 \tau T^j=T^{-j}\tau,\qquad \sigma T^j=T^{-j}\sigma .
\]

Central inversion of the plane lifts to an involution \(I\) commuting
with \(T\). On either real component specified by a fixed direction of
traversal around the inner ellipse, \(I\) and \(T\) preserve orientation.
That component is a circle. Average a metric over the finite abelian
group they generate: it acts by rotations of the resulting metric
circle, which has just one nonidentity rotation of order two.
Both \(I\) and \(T^m\) have that order and act nontrivially. They coincide
on the real circle, and therefore on connected \(X\):
\begin{equation}\label{central}
 I=T^m,\qquad P_{i+m}=-P_i .
\end{equation}
This includes primitive stars; their winding numerator is odd for even
least period. Nonidentity translations of \(X\) have no fixed point.

The quotient \(Y=X/\langle T\rangle\) is again a smooth compact genus-one
curve, and \(X\to Y\) is unramified. Cyclic area sums descend meromorphically
to \(Y\). Orientation reversal is induced by \(\tau\), with \(M\) held
fixed, so
\begin{equation}\label{reverse}
 A\circ\tau=-A,\qquad B_M\circ\tau=-B_M .
\end{equation}

\begin{lemma}\label{two}
On \(Y\), \(A\) has exactly two distinct simple poles.
If their sum is \(D\), then \(1,A\) is a basis of \(L(D)\).
\end{lemma}
\begin{proof}
The only coordinate poles of \(P\) are above the two infinite points
\([a:\pm ib:0]\) of \(E\). Neither belongs to \(C\); each has two
unramified flag preimages with simple coordinate poles. Inversion fixes
the base point and, being a nonidentity translation by \eqref{central},
exchanges its two flags. Thus their images on \(Y\) form a set \(S\)
of at most two points.

Consecutive original vertices cannot both be infinite. If distinct,
their chord is the line at infinity, not a tangent to the nondegenerate
\(C\). If they coincide, the chord is tangent to \(E\) there. Its normal
is proportional to \((1/a,\pm i/b)\), whose \(C\)-dual value is
\(\alpha^2/a^2-\beta^2/b^2=\lambda(1/b^2-1/a^2)\ne0\);
it is not tangent to \(C\). Each shoelace summand therefore has at most
one simple-pole factor. All poles of \(A\) are simple and supported on \(S\).

Choose a real orientation with \(C\) on the left of each side.
Since the origin is strictly inside \(C\),
\(\det(P_i,P_{i+1})>0\), also for a star, and \(A>0\).
By \eqref{reverse} it cannot be globally constant. A genus-one curve has
no nonconstant function with at most one simple pole. Hence \(S\) has
two distinct points, and neither possible pole of \(A\) disappears.
Riemann--Roch gives \(\dim L(D)=2\), proving the assertion.
\end{proof}

\section{The pedal trace has the same pole bound}
Write \(H=\diag(a^{-2},b^{-2})\), \(R(P)=HP\), and
\(d(P)=R(P)\cdot R(P)\), where after complexification the dot product is
bilinear, not Hermitian. The foot map is rational:
\begin{equation}\label{foot}
 Q(P;M)=M+\frac{(1-R(P)\cdot M)R(P)}{d(P)} .
\end{equation}
At \(E\)'s infinite points, \(R\) has simple leading pole direction
\((1/a,\pm i/b)\) of nonzero squared norm \(1/a^2-1/b^2\).
The quotient in \eqref{foot} is bounded, including for nonzero \(M\).

The remaining possible poles lie at four distinct finite points:
\begin{equation}\label{iso}
 d(P)=0,\qquad x^2=\frac{a^4}{c^2},\qquad y^2=-\frac{b^4}{c^2}.
\end{equation}
The zeros of \(d|_E\) are simple. The tangent \(R\cdot Z=1\) is
also tangent to \(C\), since
\[
 \alpha^2R_x^2+\beta^2R_y^2
 =a^2R_x^2+b^2R_y^2-\lambda d(P)=1 .
\]
Conversely, subtracting the two dual tangent equations forces \(d=0\).
Moreover
\[
 \frac{x^2}{\alpha^2}+\frac{y^2}{\beta^2}-1
 =-\frac{\lambda^2}{\alpha^2\beta^2}\ne0 .
\]
Thus \(\pp\) is unramified at these contacts and every individual foot
has at most a simple pole.

Let \(z\) be the common-tangent flag at one of these points; then
\(\sigma z=z\). Its other flag is \(Tz=\tau z\).
Because \(m\) is odd, choose \(k\) with \(2k\equiv1-m\pmod{2m}\).
For \(w=T^kz\),
\[
 \tau w=T^{1-k}z=T^{m+k}z=Iw .
\]
Since \(\pp\tau=\pp\), \(\pp(w)\) is fixed by planar inversion.
The only such points on projective \(E\) are its infinite points.
Every possible pedal pole is therefore in an orbit above \(S\).

It remains to exclude a double pole in the area trace. In the orbit
of a \(\sigma\)-fixed \(z\), another \(T^jz\) is \(\sigma\)-fixed exactly
when \(2j\equiv0\pmod N\), namely \(j=0,m\).
Every flag over a contact \eqref{iso} is either such a fixed flag or
\(T\) times one. The complete possible list of singular foot indices
is therefore
\[
 0,\ 1,\ m,\ m+1 .
\]
The first two original vertices tend to the same \(P\); the latter two
tend to \(-P\). Since \(m\geq3\), the only adjacent singular pairs are
\((0,1)\) and \((m,m+1)\). In a local parameter \(t\), the residues in
each pair are scalar multiples of the same vector \(R(P)\), respectively
\(R(-P)\), by \eqref{foot}. Their determinant is zero. Thus
\(\det(Q_i,Q_{i+1})\) has no \(t^{-2}\) term, even if a residue vanishes
for a special \(M\). Every other determinant has at most one singular
factor. Consequently \(B_M\) has at most simple poles above \(S\).
Unramified descent proves \(B_M\in L(D)\).

Lemma~\ref{two} gives \(B_M=c_M A+d_M\) for phase-independent constants.
Apply \eqref{reverse} and compare with the negative of this identity:
\(d_M=-d_M\), hence
\begin{equation}\label{proportional}
 B_M=c_M A .
\end{equation}
No area has been divided by. The parity congruence is essential: it
has no solution for even \(m\), and this argument asserts no converse
or other-parity classification.

\section{The credited original/outer product in the required parity}
\label{aa}
Ch\'avez-Caliz's even-period product theorem \cite[Theorems 3, 6]{chavez}
already supplies \(AA'=\text{constant}\). The following verification
under our narrower hypotheses makes its local application explicit.

Let \(\rho:X\to C^*\) remember the chord \(L\), so \(\rho\sigma=\rho\).
For \(L=[u:v:w]\), meaning \(ux+vy-wz=0\), its pole relative to \(E\) is
\[
 W=[a^2u:b^2v:w],\qquad
 \alpha^2u^2+\beta^2v^2=w^2 .
\]
It is the intersection of the \(E\)-tangents at the two chord endpoints,
so the area trace of \(W(T^i x)\) is exactly \(A'\).

A possible infinite \(W\) has \(w=0\): its line is one of the two complex
tangents to \(C\) through the origin. Its direction
\((\pm i\alpha,-\beta)\) has nonzero \(E\)-quadratic value
\[
 -\alpha^2/a^2+\beta^2/b^2=\lambda(1/a^2-1/b^2)\ne0.
\]
It meets \(E\) at two distinct finite opposite points, so \(\rho\)
is unramified there. The dual conic meets \(w=0\) transversely;
therefore \(W\) has only simple coordinate poles.
Consecutive \(W\)'s cannot both be infinite: different such chord lines
meet at the origin, which is not on \(E\); coincident successive lines
would force their intermediate endpoint to be their \(C\)-contact
at infinity, which is not on \(E\).
Thus \(A'\) has at most simple poles. Call their possible quotient
support \(R_0\).

Every orbit of \(S\) contains a \(\sigma\)-fixed flag. Indeed above
an original infinity, \(\tau x=Ix=T^m x\), so
\(\sigma x=T^{m-1}x\); solve \(2k\equiv m-1\pmod{2m}\).
Every orbit of \(R_0\) contains a \(\tau\)-fixed flag:
its two finite opposite endpoints give \(\sigma x=Ix=T^mx\),
hence \(\tau x=T^{m+1}x\); solve \(2k\equiv m+1\pmod{2m}\).
These congruences are soluble because \(m\) is odd.

The two supports are disjoint. An orbit containing \(\sigma z=z\)
has \(\tau T^jz=T^{1-j}z\); a \(\tau\)-fixed flag there would require
\(2j\equiv1\pmod N\), impossible for even \(N\).
At a \(\tau\)-fixed flag the original vertex list is palindromic,
\(\pp(T^jz)=\pp(T^{-j}z)\); its finite signed area is zero.
Thus \(A\), regular on \(R_0\), cancels every possible simple pole of \(A'\).
At a \(\sigma\)-fixed flag the same statement holds for
\(\rho(T^jz)=\rho(T^{-j}z)\) and \(W\).
Hence \(A'\), regular on \(S\), cancels every possible simple pole of \(A\).
All remaining points are regular, so \(AA'\) is holomorphic on compact
\(Y\), and constant. Multiplying \eqref{proportional} by \(A'\)
proves Theorem~\ref{thm}.

For real phases all outer intersections are finite: parallel boundary
tangents would have antipodal contacts and a center-crossing chord,
which cannot be tangent to the proper inner ellipse.
Also \(d(P)>0\) for real boundary points, so every real \(M\) is allowed.
Reversal negates both factors. Signed zero pedal areas, crossings,
and feet outside side segments create no exception.

\section{Classical reduction, exact controls, and limits}
If \(\psi_i\) is an outward unit-normal angle, central pairing gives the
classical Steiner-type identity
\begin{equation}\label{radial}
 B_M=B_0+\frac{|M|^2}{8}\sum_i\sin\bigl(2(\psi_{i+1}-\psi_i)\bigr).
\end{equation}
For completeness, in complex notation write a line as
\(\operatorname{Re}(z\overline n_i)=h_i\), \(|n_i|=1\).
Its foot is \(q_i=m/2+h_in_i-\overline m n_i^2/2\).
Expand its signed area. Translation removes \(m/2\); central pairing
changes \(h_in_i\) to its negative and leaves \(n_i^2\) unchanged,
so the linear terms cancel, leaving \eqref{radial}.
The radial fixed-polygon theorem is classical, not a discovery here;
see Steiner \cite[Sections XV--XVI, printed pp.~19--21]{steiner}
and Ferudun \cite[Section 6]{ferudun10}.
Radiality alone does not prove phase invariance: both \(A'B_0\) and
\(A'\sum_i\sin(2(\psi_{i+1}-\psi_i))\) must be constant.

As a check, when \(a=2,b=1,\lambda=4/9\), two symmetry-axis hexagons give
\[
 A'B_{(u,v)}=40\left(1+\frac{u^2+v^2}{9}\right).
\]
The general constant \(ab(a+2b)(2a+b)(1+|M|^2/(a+b)^2)\) already appears
in \cite[Proposition 4]{ferudun10}; it is not an additional new result.

Least period cannot be replaced by the number of listed entries.
For \(a=2,b=1,\lambda=4(2\sqrt{13}-5)/9\), the triangle starting at
\((2,0)\) and the triangle starting at \((0,1)\) are physical billiards
with the same caustic. Traverse each twice and set \(M=(1,0)\).
Their six-entry products, checked exactly in the supplement, are
\[
 \frac{1166+322\sqrt{13}}9,\qquad
 \frac{956+244\sqrt{13}}9.
\]
They differ and have least period three. This exemplifies the same
qualification already made in \cite[Section 7]{ferudun10}.
Odd repetitions of an eligible even primitive orbit multiply each area
by the repetition count and preserve its invariant; this does not
justify even repetitions of odd primitive orbits.

The supplied verification package contains exact reflection, caustic
tangency, general-axis hexagon and repeated-triangle controls, and an
independent high-precision physical-orbit sweep over least periods
\(6,10,14,18,22\), all tested coprime windings and several aspect ratios.
The numerical checks are finite, noninterval diagnostics, not a universal
proof. Original assertion-based checks enforce their conditions only in
ordinary Python; optimized output is not independent validation.
The later explicit-exception controls reject deliberate false identities
in both modes. Neither computations nor automated adversarial reviews
constitute formal proof-assistant verification.

The theorem excludes circular boundary ellipses, hyperbolic or degenerate
caustics, arbitrary repeated lists and unsigned filled-face areas.
The exact nonconfocal controls in the package show why a general
concentric-pair extension also requires a different argument.

\section{Chronology and disclosure}
The observed GitHub record for \href{https://github.com/AlecKriebel/Math/pull/201}
{PR 201}, with its original submitted head preserved in the supplement,
has creation time 1 October 2026, 05:18:49 UTC.
The currently exposed Zenodo record-creation times for Ferudun's v1.0
and v1.1 are respectively 1 October 2026, 16:56:14.316216 UTC and
18:12:36.529602 UTC; DataCite registration fields give 16:56:15 and
18:12:37 UTC. These establish contemporaneous separate records on that
date, not either proof's first public visibility or derivational
independence. Earlier website disclosure and other prior work are not
excluded. No first-proof, new-method, or absolute-priority claim is made.
The bounded primary-source audit and its access/coverage limits are
recorded in the supplement.

The original submission records one substantive author turn
(QUEUE \(1/5\)); the actual model and reasoning effort are unverified.
The expanded checks, source correction and this manuscript were prepared
during the October 16 audit and are not backdated.
AI tools were used extensively in the derivation, checking, literature
audit, coding and writing. This is an unrefereed research note; it has
not undergone conventional independent human peer review. Automated
adversarial audits are disclosed as such, without claiming guaranteed
correctness or historical independence.

\begin{thebibliography}{9}
\bibitem{rgk}
D. Reznik, R. Garcia and J. Koiller,
\emph{Fifty new invariants of \(N\)-periodics in the elliptic billiard},
Arnold Math. J. \textbf{7} (2021), 341--355.
\href{https://doi.org/10.1007/s40598-021-00174-y}{doi:10.1007/s40598-021-00174-y}.
\bibitem{rgkarxiv}
D. Reznik, R. Garcia and J. Koiller,
\emph{Eighty New Invariants of N-Periodics in the Elliptic Billiard},
arXiv:2004.12497v11 (2020).
\href{https://arxiv.org/abs/2004.12497v11}{arXiv record}.
Registered title; the v11 PDF has the shorter title
\emph{Eighty New Invariants in the Elliptic Billiard}.
\bibitem{chavez}
A. C. Ch\'avez-Caliz,
\emph{More about areas and centers of Poncelet polygons},
Arnold Math. J. \textbf{7} (2021), 91--106.
\href{https://doi.org/10.1007/s40598-020-00154-8}{doi:10.1007/s40598-020-00154-8}.
\bibitem{ferudun10}
A. Ferudun, \emph{Outer-Pedal Area Products in Elliptic Billiards},
unrefereed preprint, version 1.0, 1 October 2026.
\href{https://doi.org/10.5281/zenodo.23087221}{doi:10.5281/zenodo.23087221}.
\bibitem{ferudun11}
A. Ferudun, \emph{Outer-Pedal Area Products in Elliptic Billiards},
unrefereed preprint, version 1.1, 1 October 2026.
\href{https://doi.org/10.5281/zenodo.23088516}{doi:10.5281/zenodo.23088516}.
\bibitem{steiner}
J. Steiner, \emph{\"Uber den Kr\"ummungs-Schwerpunkt ebener Curven},
offprint of the memoir read 5 April 1838, printed pp.~19--21.
ETH-Bibliothek Z\"urich, Rar 5481.
\href{https://doi.org/10.3931/e-rara-3816}{doi:10.3931/e-rara-3816}.
\end{thebibliography}
\end{document}
